\chapter{Introduction}
\label{ch:intro}

\section{Host Laboratory and Dataset Description}

This report documents a second-year engineering research internship conducted from
29 June to 28 August 2026 at the \host{}, under the supervision of \supervisor{}.
The laboratory holds a large dataset of anonymized mobile network records provided by a Czech mobile
telecommunications operator, which serves as the primary dataset for several research projects conducted by the laboratory's interns.

These records contain no personally identifiable information, such as subscriber names or phone numbers, nor message content.
Each record line specifies an anonymized subscriber identifier, the temporal timestamp, and the sequence of cellular antennas to which the mobile terminal was attached over a 24-hour period. This anonymized corpus constitutes the empirical foundation of the work described in this report.

\section{Research Objectives and Next-Cell Prediction Task}

Mobile terminals continuously attach to the cellular antenna offering the highest signal quality. As users move, the terminal handsover connection from one antenna to the next. The mobile operator logs an event line upon each cell transition, as well as periodically during stationary periods.

The central research objective of this internship is formulated as follows: \emph{given the initial segment of a subscriber's daily mobility trajectory, can the model accurately predict the next cellular antenna to which the device will connect?} Predicting cellular transitions enables mobile network operators to optimize resource allocation, anticipate local capacity demands, and reduce connection drop rates through proactive handover initialization.

A fundamental distinction must be established between physical infrastructure and logical coverage areas.\footnote{\textbf{Base station} (or \emph{site}/\emph{mast}): a physical infrastructure site. \textbf{Cell}: the spatial coverage area served by an individual directional antenna mounted on a base station.} A \textbf{base station} corresponds to a physical mast housing multiple directional antennas. Each individual antenna defines a distinct \textbf{cell}. The geographical study area evaluated in this work comprises \textbf{369 cells} distributed across \textbf{113 base stations}, representing an average of three to four cells per physical site. The target of the predictive task is the specific cell rather than the base station, which presents a higher predictive complexity than coarse-grained spatial localization.

\section{Next-Cell Prediction via Fine-Tuned Language Models}

As illustrated in \textbf{Figure~\ref{fig:overview}}, the overall experimental pipeline processes raw network logs into a ranked list of candidate cellular antennas, evaluated against a reference baseline.

% F0 -- The whole project on one line, before any vocabulary is introduced.
\begin{figure}[!htbp]
\centering
\fitwidth{%
\begin{tikzpicture}[font=\small,
  st/.style={draw=clrule,fill=clsoft,rounded corners=3pt,align=center,
             inner sep=7pt,text width=3.15cm,minimum height=3.55cm,anchor=north},
  ar/.style={-{Stealth[length=6pt]},clrule,very thick},
  num/.style={circle,draw=clmodel,fill=white,inner sep=1.4pt,
              font=\scriptsize\bfseries,text=clmodel},
]
\node[st] (s1) at (0,0)
  {{\footnotesize\bfseries\color{clmodel} 1. One day of\\ one person}\\[4pt]\footnotesize
   The network writes down which antenna the phone was on, and at what time,
   all day long.};
\node[st] (s2) at (4.35,0)
  {{\footnotesize\bfseries\color{clmodel} 2. Rewritten\\ as a list}\\[4pt]\footnotesize
   That day becomes an ordered list of \emph{antenna, hour} pairs, one line
   per person.};
\node[st] (s3) at (8.70,0)
  {{\footnotesize\bfseries\color{clmodel} 3. A language\\ model reads it}\\[4pt]\footnotesize
   The same kind of model that finishes a sentence, retrained on those lists
   instead of on text.};
\node[st] (s4) at (13.05,0)
  {{\footnotesize\bfseries\color{clmodel} 4. It names the\\ next antenna}\\[4pt]\footnotesize
   Out of the 369 antennas of the study area, it ranks the likeliest one
   first.};

\draw[ar] (s1.east) -- (s2.west);
\draw[ar] (s2.east) -- (s3.west);
\draw[ar] (s3.east) -- (s4.west);
\node[num] at (s1.north west) {1}; \node[num] at (s2.north west) {2};
\node[num] at (s3.north west) {3}; \node[num] at (s4.north west) {4};

% --- the competitor, underneath -----------------------------------------
\node[draw=clbase,dashed,fill=clbase!8,rounded corners=3pt,align=center,
      inner sep=7pt,text width=14.1cm,font=\small,anchor=north]
      (rule) at (6.52,-4.25)
  {\textbf{The thing to beat, on the very same question:} answer
   \emph{``the same antenna as a moment ago''}, every time. No model, no
   training, one line of code.\\[2pt]
   On an ordinary day of records it is right \textbf{69 times out of 100}.};
\draw[ar,clbase] (s4.south) -- (13.05,-4.25);
\end{tikzpicture}}
\caption{What this internship built, in four steps}
\label{fig:overview}
\figtext{%
\figpara{What the picture says.}{Reading left to right: the mobile network
records, all day, which antenna each phone is attached to. I turn one person's
day into a plain ordered list. A small language model, the same kind of program
that predicts the next word of a sentence, is retrained on thousands of those
lists. Asked to continue one of them, it names which antenna the person will be
on next, and it does so by ranking all 369 antennas of the area from the most
to the least likely.}

\figpara{Why a language model, for a problem with no language in it.}{Because
the shape of the problem is the same. A sentence is an ordered list of words
drawn from a fixed vocabulary, and the model's job is to guess the next one. A
day of movement is an ordered list of antennas drawn from a fixed set of 369,
and the job is exactly the same. Nothing else about language is used here.}

\figpara{The dashed box is the real difficulty.}{A competitor with no model at
all, ``the person is still where they were'', is already right about seven times
out of ten, because people mostly stay put. Every number in this report is the
distance between my model and that rule, measured on the same people, on the
same moments of their day.}}
\end{figure}


At the core of this architecture is a \textbf{language model}.\footnote{\textbf{Language model}: a probabilistic model trained to estimate the probability distribution over sequences of discrete tokens. The underlying mechanics operate on arbitrary discrete sequences, allowing its application to spatial mobility traces.} While language models are conventionally employed for natural language processing, their mathematical framework operates on generic ordered sequences of discrete tokens. In this context, daily mobility traces are framed as ordered sequences of cell identifiers selected from a fixed vocabulary of 369 discrete cells.

Rather than using off-the-shelf general language models directly, parameter-efficient \textbf{fine-tuning} (PEFT) is applied. The internal parameters of the neural network are adapted using local mobility traces, specializing the model from a general natural language generator into a dedicated next-cell predictor for the target geographical region, as detailed in Chapter~\ref{ch:tools}. Alternative approaches explored in parallel research rely on prompting off-the-shelf models without parameter updates, addressing a distinct research scope.

\section{Baseline Evaluation: The Persistence Baseline}

To rigorously assess model performance, predictions must be evaluated against a strong deterministic benchmark: the \textbf{persistence baseline}.

The persistence baseline evaluates the heuristic rule: \emph{``the next cell is identical to the current cell''}. Despite its simplicity, this single-line baseline achieves an accuracy of \textbf{69.2\%} on raw daily network records, primarily because individuals remain stationary for large portions of the day. Beating this persistence baseline constitutes the key criterion for demonstrating true predictive capability. Across early experimental iterations, baseline performance proved difficult to surpass.

% F1 -- What a single prediction is. Deliberately the simplest figure in the
% report: three columns A, B, C, and two rows, one per kind of moment.
\begin{figure}[!htbp]
\centering
\fitwidth{%
\begin{tikzpicture}[font=\small,
  ev/.style={draw=clrule,fill=clsoft,rounded corners=2pt,inner sep=3pt,
             font=\scriptsize\ttfamily,minimum width=1.35cm,minimum height=0.55cm},
  tgt/.style={draw=clmodel,fill=white,thick,rounded corners=2pt,inner sep=3pt,
              minimum width=1.35cm,minimum height=0.55cm},
  hr/.style={font=\scriptsize,text=clrule},
  hdr/.style={font=\footnotesize\bfseries,text=clmodel,anchor=west},
  cs/.style={font=\footnotesize\bfseries,anchor=east},
  ans/.style={font=\scriptsize,anchor=west},
  note/.style={font=\scriptsize,text=clrule,anchor=west},
]
% ================= column headers =================
\node[hdr] at (-0.70,2.05) {A. What the model is given};
\node[hdr] at (5.55,2.05)  {B. The moment};
\node[hdr] at (8.30,2.05)  {C. The two answers, compared};
\draw[clrule] (-0.80,1.80) -- (14.2,1.80);

% ================= row 1: the person has not moved =================
\node[cs] at (-1.00,0.85) {1};
\foreach \i/\h/\c in {0/07h/BKVCHE1, 1/09h/BKVCHE1, 2/11h/BKVCHE1, 3/12h/BKVCHE1}{
  \node[hr] at (\i*1.50,1.42) {\h};
  \node[ev] at (\i*1.50,0.95) {\c};
}
\node[hr]  at (6.15,1.42) {13h};
\node[tgt] at (6.15,0.95) {\textbf{?}};
\node[ans] at (8.30,1.32) {really visited: \cid{BKVCHE1}};
\node[ans] at (8.30,0.97) {the rule says \cid{BKVCHE1}: \textcolor{clgain}{\textbf{right}}};
\node[ans] at (8.30,0.62) {my model says \cid{BKVCHE1}: \textcolor{clgain}{\textbf{right}}};
\node[note] at (-0.70,0.20)
  {\textbf{The person has not moved.} Both answers are free, and neither is
   about the model. This is the ordinary case.};

\draw[clrule,dotted] (-0.80,-0.18) -- (14.2,-0.18);

% ================= row 2: the person has gone somewhere =================
\node[cs] at (-1.00,-1.15) {2};
\foreach \i/\h/\c in {0/07h/BKVCHE1, 1/09h/BKVCHE1, 2/11h/BKVCHE1, 3/12h/BKVCHE1}{
  \node[hr] at (\i*1.50,-0.58) {\h};
  \node[ev] at (\i*1.50,-1.05) {\c};
}
\node[hr]  at (6.15,-0.58) {18h};
\node[tgt] at (6.15,-1.05) {\textbf{?}};
\node[ans] at (8.30,-0.68) {really visited: \cid{BKVPAN1}};
\node[ans] at (8.30,-1.03) {the rule says \cid{BKVCHE1}: \textcolor{clloss}{\textbf{wrong}}};
\node[ans] at (8.30,-1.38) {my model says \cid{BKVPAN1}: \textcolor{clgain}{\textbf{right}}};
\node[note] at (-0.70,-1.80)
  {\textbf{The person has gone somewhere.} The rule must fail. This is the only
   kind of moment on which I can win anything.};
\end{tikzpicture}}
\caption{What one prediction is}
\label{fig:worked-example}
\figtext{%
\figpara{How to read it.}{The two rows are two different moments of the same
person's day, and the three columns are always the same three things. Column A
is the past the model is allowed to read: four moments, each written as an hour
and the name of the antenna the phone was on. Column B is the moment I want an
answer for, with the hour given and the antenna hidden. Column C compares the
answer of the one-line rule with the answer of my model against the antenna the
person really visited.}

\figpara{Row 1: nothing happens.}{The person is still on the same antenna at
13h. Answering ``the same as before'' is correct, so the one-line rule scores a
point without knowing anything, and my model scores the same point. Moments like
this one are the majority, which is why a rule with no model in it is right
about seven times out of ten.}

\figpara{Row 2: the person leaves.}{At 18h they are on a different antenna. The
rule now has to be wrong, because being wrong is all it can do when someone
moves. This is the only kind of moment on which a model can gain anything at
all, and every gain reported in this document was won on moments of this second
kind. The hour is what makes the second row answerable: 18h is when this
particular person leaves.}

\figpara{One caution.}{The two rows are drawn to explain the task, not measured
from a run. The antenna names are real names from the study area; how often each
of the two cases actually occurs is measured in Chapter~\ref{ch:data}.}}
\end{figure}


\textbf{Figure~\ref{fig:worked-example}} illustrates the mechanics of the persistence baseline compared to the model predictions across two subscriber trajectories. When a user is stationary, repeating the previous cell identifier yields a correct prediction for both the baseline and the model in \textbf{Figure~\ref{fig:worked-example}}. However, when a transition occurs, the persistence baseline systematically fails, requiring non-trivial trajectory modeling. Performance improvements reported in this document are measured exclusively on non-stationary transition events.

Throughout this report, performance gains are reported as the \emph{difference in accuracy percentage points} relative to the persistence baseline. A reported gain of $+4.3\%$ indicates that the fine-tuned model achieves an accuracy of 54.5\% on a benchmark where the persistence baseline yields 50.2\% across identical evaluation instances. Evaluating relative performance delta ensures consistent comparisons across dataset revisions and experimental cohorts.

\section{Related Work and Collaborative Context}

This project was conducted alongside research by peer group members investigating the same mobile network dataset. \textbf{Figure~\ref{fig:team-context}} illustrates the placement of this study within the broader team research framework.

% F1 -- Where this internship sits inside the team's work on the same corpus.
\begin{figure}[t]
\centering
\fitwidth{%
\begin{tikzpicture}[
  font=\small,
  box/.style={draw=clrule,fill=clsoft,rounded corners=2pt,align=center,
              inner sep=5pt,minimum height=1.05cm,text width=3.5cm},
  mine/.style={box,draw=clmodel,fill=clmodel!8,text width=3.9cm},
  corpus/.style={box,draw=clrule,fill=white,text width=11.6cm},
  yard/.style={draw=clbase,dashed,fill=clbase!8,rounded corners=2pt,align=center,
               inner sep=5pt,text width=11.6cm},
  ar/.style={-{Stealth[length=4pt]},clrule,thick},
]
\node[corpus] (raw) {\textbf{Shared corpus} --- 15 daily files, 12--26 March 2014,
  Karlovy Vary study area\\[1pt]
  \footnotesize 80{,}217 to 99{,}504 users per day $\cdot$ 369 cells over 113 base stations
  $\cdot$ one row = one user, one day};

\node[box,below left=1.0cm and 0.05cm of raw.south] (nino) {\textbf{Nino}\\[1pt]
  \footnotesize cleaning, per-day exports, mobile-user selection};
\node[box,right=0.35cm of nino] (tb) {\textbf{Tiphain \& Benjamin}\\[1pt]
  \footnotesize move-or-stay classification, prompting a frozen model};
\node[box,right=0.35cm of tb] (tip) {\textbf{Tiphain}\\[1pt]
  \footnotesize equal record count per user (200 records)};

\node[mine,below=1.15cm of tb] (me) {\textbf{This internship}\\[1pt]
  \footnotesize \emph{fine-tune} a language model to predict a user's next \cid{cell\_id}};

\node[yard,below=0.75cm of me] (yardn) {\textbf{The common yardstick: the persistence baseline}
  --- ``predict the last-seen \cid{cell\_id}'', no model at all.\\[1pt]
  \footnotesize Every claim in this report is stated as a gap to this baseline, measured on the very same data.};

\draw[ar] (raw.south) ++(-3.9,0) -- (nino.north);
\draw[ar] (raw.south) -- (tb.north);
\draw[ar] (raw.south) ++(3.9,0) -- (tip.north);
\draw[ar] (nino.south) -- ++(0,-0.35) -| (me.north west);
\draw[ar] (tb.south) -- (me.north);
\draw[ar] (tip.south) -- ++(0,-0.35) -| (me.north east);
\draw[ar] (me.south) -- (yardn.north);
\end{tikzpicture}}
\caption[Where this internship sits inside the team's work]{Where this
internship sits. Three colleagues worked on the same 15-day corpus with
different instruments; I inherited their cleaned exports and their reusable
ideas, and contributed the fine-tuning branch. Everything below the dashed
line is what makes the branches comparable: a single, model-free yardstick.}
\label{fig:team-context}
\end{figure}


Collaborative efforts involved trajectory preprocessing and dataset characterization. Nino developed corpus extraction scripts for identifying user home/activity cells and filtering active mobility cohorts, providing the daily trajectory extractions utilized in this work. Tiphain and Benjamin investigated zero-shot prompting techniques on general LLMs, while Reda investigated complementary mobility pattern analysis.

\section{Report Structure}

This report is organized by thematic research questions rather than chronological execution; the complete timeline of experimental iterations is provided in Appendix~\ref{app:timeline}. Chapters~\ref{ch:data} through \ref{ch:scale} conclude with a synthesis section summarizing empirical conclusions. Appendix~\ref{app:repro} details experimental hyperparameters and tabular data provenance, while Appendix~\ref{app:glossary} provides a reference glossary of technical domain terminology.


